% Generated by member4_training_eval/make_results_text.py -- do not edit by hand.
\subsection{Main Results}
\label{sec:main_results}

Table~\ref{tab:main_results} gives the accuracy of every group on the in-domain Photo test set and on the three unseen target domains, averaged over 5 seeds.

The Original group reaches $40.8\pm2.3\%$ on the cross-domain average while scoring above 99\% on Photo Test, so the domain gap is large and unambiguous. Adding the AdaIN-styled images raises the cross-domain average to $47.4\pm0.8\%$, a paired gain of $+6.53$ points (95\% CI $[+3.13, +9.93]$, $t=+5.33$, $p=0.0060$, significant over the shared seeds).

\textit{Where the gain comes from.} The styled groups have twice the training rows of the Original group, so the headline difference cannot be attributed to style until the data volume is controlled. The two control groups do exactly that. 

+Duplicate, which adds a byte-identical copy of every training image, reaches $37.7\%$: the pure effect of doubling the training rows is $-3.14$ points (95\% CI $[-5.81, -0.46]$, $p=0.0312$, significant). 
+NonStyle, which adds a non-stylistic transform instead, reaches $39.9\%$, a further $+2.19$ points over +Duplicate (95\% CI $[+0.04, +4.33]$, $p=0.0473$, significant). 
Style on top of that generic appearance change contributes $+7.48$ points (95\% CI $[+6.09, +8.87]$, $p=0.0001$, significant). 

The third row of Table~\ref{tab:decomposition} is an \emph{upper bound}: it contains the comparison of styling against a generic transform and the comparison of one particular family of transforms against a random one at the same time, so it is not a strict causal attribution of style.

\textit{The spatial modulation adds nothing.} The strength-matched control, which receives exactly the same average style strength as +Ours but distributes it uniformly, reaches $45.4\%$, against $45.2\%$ for +Ours: a paired difference of $-0.17$ points (95\% CI $[-1.71, +1.37]$, $t=-0.31$, $p=0.7750$, not significant). 
The equal-strength spatial map confirms the same thing from the other side: with the mean strength pinned to 0.500 it reaches $47.2\%$ against +AdaIN's $47.4\%$, a paired difference of $-0.15$ points (95\% CI $[-1.42, +1.11]$, $p=0.7563$, not significant). 

Once the average strength is held fixed, the two methods are indistinguishable: the accuracy of a styled training set is governed by \emph{how much} style it applies, not by \emph{where} the style is placed.
